\section{Точные конструкции в размерностях $4$, $5$, $7$ и $9$}
\label{sec:constr}

Каждая из четырёх теорем этого раздела предъявляет рациональную решётку и
подрешётку и доказывается протоколом предложения~\ref{prop:protocol}; ниже
приводятся данные конструкции и величины, которые протокол предъявляет.
Полные списки вершин, векторов окна и KKT-сертификатов лежат в файлах,
указанных в таблице~\ref{tab:manifest}. О том, как конструкции были найдены,
сказано по одному абзацу; подробности поиска --- в~\cite{Repo}.

\subsection{Размерность $4$: $\chi(\R^4)\le43$}\label{ssec:dim4}

\paragraph{Эйзенштейново семейство.}
Пусть $S\in\mathrm{GL}_4(\Z)$ --- элемент порядка~$3$ с характеристическим
многочленом $\Phi_3^2$ и $\Lambda$ --- решётка с $S\Lambda=\Lambda$. Тогда
$\Lambda$ --- $\Z[\omega]$-модуль ранга~$2$ ($\omega=e^{2\pi i/3}$), а
$S$-инвариантные положительно определённые формы образуют
четырёхпараметрическое семейство, определённое над~$\Q$: эрмитова форма
$h=\bigl(\begin{smallmatrix}a&c\\ \bar c&b\end{smallmatrix}\bigr)$ задаёт
вещественное произведение $\langle x,y\rangle=\operatorname{Re}h(x,y)$, и в
$\Z$-базисе $(e_1,\omega e_1,e_2,\omega e_2)$ её матрица Грама при
$u=\operatorname{Re}c$, $s=\tfrac{\sqrt3}{2}\operatorname{Im}c$ равна
\begin{equation}\label{eq:eisform}
Q(a,b,u,s)=\begin{pmatrix}
a & -\tfrac a2 & u & -\tfrac u2-s\\[2pt]
-\tfrac a2 & a & -\tfrac u2+s & u\\[2pt]
u & -\tfrac u2+s & b & -\tfrac b2\\[2pt]
-\tfrac u2-s & u & -\tfrac b2 & b
\end{pmatrix}.
\end{equation}
Если и подрешётка $S$-инвариантна, то она --- $\Z[\omega]$-подмодуль, и её
индекс есть норма идеала; $43=N(7+\omega)$ --- норма эйзенштейнова простого,
и подмодулей индекса~$43$ ровно $2\cdot44=88$. Сужение поиска до этого
семейства (три параметра вместо девяти, $88$ подрешёток вместо
$1+43+43^2+43^3=81\,400$) и дало конструкцию ниже; в полном пространстве форм
поиск останавливался ниже порога~\cite{Repo}.

\paragraph{Конструкция.}
Пусть $\Lambda_\star=B^{\mathsf T}\Z^4$, где $Q=BB^{\mathsf T}$ --- форма
\eqref{eq:eisform} с параметрами
\begin{equation}\label{eq:gram43}
a=1,\qquad b=\frac{39261}{25000},\qquad u=\frac{37381}{100000},\qquad
s=\frac{3141}{12500}
\qquad(200\,000\,Q\in\mathrm{Mat}_4(\Z)),
\end{equation}
и пусть $\Gamma_\star$ --- подрешётка индекса $43$ с матрицей перехода
(эрмитова нормальная форма; строки --- коэффициенты по базису $B$)
\begin{equation}\label{eq:hnf43}
T=\begin{pmatrix}1&0&0&41\\0&1&0&14\\0&0&1&37\\0&0&0&43\end{pmatrix},
\qquad\det T=43 .
\end{equation}
Инварианты (все рациональны, так как рациональна $Q$):
\[
\begin{aligned}
\lambda_1(\Lambda_\star)^2&=1,\\
\diam^2V_0&=\frac{5374343884337}{2019775849050}
  =2{,}660861544\ldots,\\
D(\Gamma_\star)^2&=\frac{7396075258066147}{2756867811550000}
  =2{,}682781969\ldots,
\end{aligned}
\qquad
\begin{aligned}
d(\Lambda_\star,\Gamma_\star)^2
&=\frac{298768283679964998359742207}{296327113258585430253847000}\\
&=1{,}008238093\ldots,\\
d&=1{,}004110598\ldots
\end{aligned}
\]
Ячейка $V_0$ имеет максимально возможные $30$ фасет и $f$-вектор
$(120,240,150,30)$. Решётка \emph{не} общего положения:
$\mathrm{Aut}(\Lambda_\star)\cong\Z/6$, три пары кратчайших векторов, и вся
группа автоморфизмов сохраняет $\Gamma_\star$, поэтому минимум $D$
достигается сразу на трёх парах $\pm v$ одной орбиты $\Z/3$.

\begin{keybox}
\begin{theorem}\label{thm:main43}
$\chi\bigl(\R^4,[1,\ell]\bigr)\le43$ при всех $1\le\ell\le\ell_0=1{,}00411$;
в частности, $\chi(\R^4)\le43$.
\end{theorem}
\end{keybox}

\begin{proof}[Доказательство --- протокол предложения~\ref{prop:protocol}]
\emph{Индекс:} $\det T=43$. \emph{Фасеты:} релевантные векторы --- минимумы
нормы в $15$ ненулевых классах $\Lambda_\star/2\Lambda_\star$, каждый класс
даёт единственную пару $\pm v$, итого $30$ фасет. \emph{Вершины:} решены все
$\binom{30}{4}=27\,405$ рациональных систем $4\times4$ с проверкой всех
неравенств; вершин ровно $120$, множество центрально симметрично,
$\vol V_0=\det\Lambda_\star$; отсюда $\diam^2V_0$ выше.
\emph{Окно}~\eqref{eq:window} при $\ell_0=1{,}00411$ перечислено полностью;
\emph{разделение:} для каждого вектора окна расстояние $\dist(v/2,V_0)$
вычислено точно --- активное множество берётся из плавающей точки, а
оптимальность подтверждается условиями Каруша--Куна--Таккера в дробях;
минимум $D^2$ выписан выше. \emph{Запас:}
\[
D(\Gamma_\star)^2-\ell_0^2\,\diam^2V_0
=\frac{3559631641205182621936510913}{1113651004958403329305500000000000}
=3{,}196\cdot10^{-6}>0 .
\qedhere
\]
\end{proof}

Независимой проверкой служит второй, \emph{опорный} сертификат: ячейка
заменяется объемлющим многогранником $P\supseteq V_0$ --- пересечением тех же
$30$ полупространств, --- для которого $\diam V_0\le\diam P$ и
$\dist(v/2,V_0)\ge\dist(v/2,P)\ge\bigl(\langle v/2,u\rangle-\max_{w\in
\mathrm{vert}P}\langle w,u\rangle\bigr)/\|u\|$ с рациональным направлением
$u$. Он даёт тот же диаметр (как дробь) и оценку снизу
$D(\Gamma_\star)^2\ge869307639027355969/324032161065600000=2{,}682781969\ldots$,
не превосходящую точного значения и превосходящую порог
$\ell_0^2\diam^2P=2{,}682778773\ldots$ Дополнительно исчерпывающим перебором
всех $81\,400$ подрешёток индекса~$43$ подтверждено, что~\eqref{eq:hnf43}
доставляет максимум $d$ для $\Lambda_\star$.

\begin{remark}[оптимум семейства]
Точка~\eqref{eq:gram43} --- рационализация (со знаменателем $10^5$) численного
оптимума задачи $\max d$ по трёхпараметрическому семейству; сам оптимум
устроен как жёсткое равноколебание (четыре $\Z/3$-орбиты векторов
$\Gamma_\star$ дают одно значение $D$, диаметр реализуется на $15$ парах
вершин) и численно равен $d=1{,}00411205721\ldots$ Исторически первой была
найдена другая решётка того же семейства с $200\,Q\in\mathrm{Mat}_4(\Z)$ и
$d^2=14984024/14873313$, $d=1{,}0037149\ldots$
(сертификат $\ell_0=1{,}003714$); именно её симметрия подсказала
семейство~\eqref{eq:eisform}.
\end{remark}

\subsection{Размерность $5$: $\chi(\R^5)\le132$}\label{ssec:dim5}

Прежняя оценка $140$ получена АБПР на решётке $A_5^*$, и там же доказана её
минимальность --- но лишь среди подрешёток \emph{этой} решётки. Конструкция
ниже найдена чередованием дискретного поиска ядра (подрешётка задаётся
сюръективным гомоморфизмом $\varphi\colon\Z^5\to\Z/132$) с непрерывной
деформацией формы Грама~\cite{Repo}; проверяется её
рационализация.

\begin{keybox}
\begin{theorem}\label{thm:r5}
$\chi\bigl(\R^5,[1,\ell]\bigr)\le132$ при всех $1\le\ell\le\dfrac{101}{100}$;
в частности, $\chi(\R^5)\le132$.
\end{theorem}
\end{keybox}

\begin{proof}[Доказательство --- протокол предложения~\ref{prop:protocol}]
Пусть $\Lambda=B^{\mathsf T}\Z^5$, где $BB^{\mathsf T}=Q$ и
\[
Q=\frac1{100000}
\begin{pmatrix}
79327&86284&69597&36140&26501\\
86284&254538&216808&101257&102135\\
69597&216808&296568&163667&140527\\
36140&101257&163667&181733&71653\\
26501&102135&140527&71653&148076
\end{pmatrix}.
\]
Её главные угловые миноры
\[
79327,\ 12746807270,\ 1422467480626358,\
124437401434750889345,\ 9999935596575703407615413
\]
положительны. Зададим цвет координатного вектора $x\in\Z^5$ тройкой
гомоморфизмов
\[
\begin{aligned}
\varphi_{11}(x)&=x_2+6x_3+2x_4+3x_5\pmod {11},\\
\varphi_3(x)&=x_1+2x_2+x_3\pmod 3,\\
\varphi_4(x)&=x_2+3x_3+x_4+3x_5\pmod 4.
\end{aligned}
\]
По китайской теореме об остатках это эквивалентно одной строке
\[
\varphi(x)=88x_1+89x_2+127x_3+57x_4+3x_5\pmod {132}.
\]
Положим $\Gamma=B^{\mathsf T}\ker\varphi$. Гомоморфизм сюръективен,
$[\Lambda:\Gamma]=132$, а диагональ Смита ядра равна $(1,1,1,1,132)$.

Ячейка Вороного рациональной формы $Q$ имеет $62$ фасеты и $720$ вершин.
Все вершинные системы решены над $\mathbb Q$; точный квадрат радиуса покрытия
равен
\[
R^2=
\frac{3093570095915736710467707313641}
     {3999974238630281363046165200000}.
\]
Доказываемое утверждение --- это $D(v)\ge\ell\,\diam V_0$ при $\ell=101/100$,
поэтому возможный нарушитель обязан удовлетворять
\begin{equation}\label{eq:win5}
\|v\|<2(1+\ell)R,\qquad\text{т.е.}\qquad \|v\|^2<12{,}4984\ldots
\end{equation}
Перебор Финке--Похста \cite{FinckePohst} по LLL-приведённому \cite{LLL}
базису $\Gamma$ оставляет в окне~\eqref{eq:win5} ровно $38$ ненулевых
векторов ($19$ пар); из них $36$ лежат уже в окне $\|v\|<4R$ и независимо
воспроизводятся C\raise.1ex\hbox{\small++}-перечислителем, а оставшаяся пара
$\pm(-2,2,-2,2,2)$ даёт $D/\diam V_0=1{,}0442\ldots$ и на минимум не влияет.

Для всех $38$ векторов задача проекции на $V_0$ решена точно: допустимость,
неотрицательность множителей KKT и значения целей проверены над $\mathbb Q$.
Минимум достигается на паре $\pm(2,0,1,-3,0)$ (координаты по базису
$\Lambda$) и равен
\[
\begin{aligned}
D_{\min}^2&=
\frac{1634114770527727}{516898614620000},\\
D_{\min}^2-\left(\frac{101}{100}\right)^2 4R^2
&=\frac{1450500657237822255902962097875907780124229}
{258447642807960215367421006872956903000000000}>0.
\end{aligned}
\]
Следовательно, $D_{\min}/\diam V_0=1{,}010897714548927\ldots>1{,}01$, и
раскраска ячеек по классам $\Lambda/\Gamma$ даёт требуемые $132$ цвета.
\end{proof}

Дополнительный верификатор, не использующий ни Qhull, ни вещественный
перечислитель коротких векторов, точно проверяет каждый из $31$ ненулевого
класса чётности, строит $62$ фасеты, обходит связный граф из $720$ вершин и
$1800$ рёбер и воспроизводит те же $R^2$, $D_{\min}^2$ и KKT-сертификаты.
Рационализация той же численной формы со знаменателями $10^4$ и $10^6$ даёт
отношения $1{,}010770760\ldots$ и $1{,}010904641\ldots$ --- в обоих случаях
строгий запас для $\ell=1{,}01$ сохраняется.

\subsection{Размерность $7$: $\chi(\R^7)\le1029$}\label{ssec:dim7}

\paragraph{Ламинирование.}
Базой служит эйзенштейнова раскраска $E_6^*/343$: $\Gamma'=(3+\omega)E_6^*$,
индекс $7^3=343$, ширина $\sqrt{7/6}$ (предложение~\ref{prop:planar} при
радиусе покрытия $E_6^*$, равном $\lambda_1/\sqrt2$~\cite{ConwaySloane}).
\emph{Ламинирование} строит раскраску $\R^7$ из раскраски $\R^6$: для
сдвига $c\in\R^6$, высоты $t>0$, склейки $a\in\Lambda'/\Gamma'$ и модуля
слоя $m$ положим
\[
\Lambda=\bigl\{(x+ic,\;it)\ :\ x\in\Lambda',\ i\in\Z\bigr\},\qquad
\Gamma=\bigl\{(x+ic,\;it)\ :\ \psi(x)+ia=0,\ m\mid i\bigr\},
\]
где $\Gamma'=\ker\psi$; тогда $[\Lambda:\Gamma]=343\,m$. При $m=3$ индекс
равен $1029$. Ячейка при подъёме растёт не более чем на $t^2$ по квадрату
диаметра, а расстояния между <<горизонтальными>> одноцветными ячейками не
убывают, так что запас ширины базы расходуется на высоту слоя; сдвиг $c$,
смещённый с орбит глубоких дыр $E_6^*$, и высота $t\approx1{,}08$ найдены
оптимизацией~\cite{Repo}. Проверяется рациональная точка: в
подходящем $\Z$-базисе семимерной решётки её матрица Грама и подрешётка
таковы:
% Сгенерировано из audit-data/results/dim7_1029_exact.json (не править вручную)
\begin{equation}\label{eq:gram1029}
\begingroup
\setlength{\arraycolsep}{3.2pt}
G=\begin{pmatrix}
30000&15000&0&-22500&15000&7500&-8632\\
15000&30000&22500&0&7500&15000&-4197\\
0&22500&45000&22500&0&22500&1601\\
-22500&0&22500&45000&-22500&0&6617\\
15000&7500&0&-22500&30000&15000&-5275\\
7500&15000&22500&0&15000&30000&4434\\
-8632&-4197&1601&6617&-5275&4434&17897
\end{pmatrix},\qquad Q=\frac{G}{10^{4}},
\endgroup
\end{equation}
\begin{equation}\label{eq:sub1029}
\begingroup
\setlength{\arraycolsep}{3.6pt}
C=\begin{pmatrix}
7&-5&0&0&0&0&0\\
0&1&0&0&0&0&0\\
0&0&7&-5&0&0&0\\
0&0&0&1&0&0&0\\
0&0&0&0&7&-5&0\\
0&0&0&0&0&1&0\\
0&0&0&0&0&0&3
\end{pmatrix},\qquad \det C=7^3\cdot3=1029 .
\endgroup
\end{equation}

Матрица $G$ имеет положительные главные угловые миноры; первые шесть базисных
векторов порождают базу, седьмой --- слоевой вектор $(c,t)$; три блока
$\bigl(\begin{smallmatrix}7&-5\\0&1\end{smallmatrix}\bigr)$ матрицы $C$
задают в трёх комплексных координатах базы идеал нормы~$7$ (умножение на
$3+\omega$ при подходящем выборе $\Z[\omega]$-структуры), множитель $3$ ---
модуль слоя; в частности, $D(\Gamma)$ достигается на образах минимальных
векторов базы.

\begin{keybox}
\begin{theorem}[точный сертификат в $\R^7$]\label{thm:r7ex}
Пусть $\Lambda=\Z^7$ со скалярным произведением $x^{\mathsf T}Qy$,
$Q=G/10^4$ из~\eqref{eq:gram1029}, и $\Gamma=C\Z^7$ с $C$
из~\eqref{eq:sub1029}, $\det C=1029$. Тогда
\[
D(\Gamma)^2=7,\qquad
\diam^2 V_0=\frac{96858928789031597}{14761819462500000},\qquad
d^2=\frac{103332736237500000}{96858928789031597},
\]
то есть $d=1{,}032878\ldots$, и $\chi\bigl(\R^7,[1,\ell]\bigr)\le1029$ при
всех $1\le\ell\le\dfrac{103}{100}$; в частности, $\chi(\R^7)\le1029$.
\end{theorem}
\end{keybox}

\begin{proof}[Доказательство --- протокол предложения~\ref{prop:protocol}]
\emph{Ячейка.} Релевантные векторы находятся по теореме Вороного о
$\Lambda/2\Lambda$; граница перебора $\max_{s\in\{0,1\}^7}Q(s,s)$ полна,
поскольку у каждого класса есть представитель из нулей и единиц; из
$300\,368$ перебранных векторов получаются $127$ пар, то есть $254$ фасеты,
максимум, возможный в $\R^7$ (ни один класс не вырождается).

\emph{Вершины.} Qhull лишь \emph{предлагает} точки; каждая переводится в
точную рациональную вершину решением целочисленной системы по $7$
независимым активным фасетам, после чего полный активный набор
пересчитывается над $\Z$. Полнота доказывается точным замыканием по
$1$-скелету: у каждой вершины перебираются все $6$-подмножества активных
фасет, и если подмножество ранга~$6$ задаёт допустимый экстремальный луч
касательного конуса, то второй конец соответствующего ребра обязан лежать в
наборе; связность $1$-скелета выпуклого многогранника доказывает полноту.
Простота \emph{не} предполагается --- и это существенно: из $30\,368$ вершин
$1\,600$ имеют по девять активных фасет. Отсюда точное
$R^2=96858928789031597/59047277850000000$ и $\diam^2V_0=4R^2$.

\emph{Разделение.} Улучшить текущий рекорд $D_*$ может лишь вектор с
$\|v\|<D_{*}+\diam V_0$; чтобы остаться в рациональных числах, берётся
подстраховка $(a+b)^2\le2(a^2+b^2)$, то есть перебор
$\|v\|^2\le2\bigl(D_{*}^2+\diam^2 V_0\bigr)=
200191665026531597/7380909731250000$. В это окно попадают $74$ вектора
($37$ пар); для каждого $D(v)^2=4\dist(v/2,V_0)^2$ предъявляется точный
KKT-сертификат --- множители $\mu\ge0$ и допустимость проекции.

\emph{Минимум} равен $7$ ровно и достигается на $27$ парах $\pm v$
\emph{горизонтальных} векторов (нулевая слоевая координата) --- образах
$54$ минимальных векторов $E_6^*$: это в точности пол
$\sqrt{7/3}\,\lambda_1$ предложения~\ref{prop:planar} в нормировке
$\lambda_1(E_6^*)^2=3$. Ближайший слоевой вектор даёт
$D^2=478385948533/68336903600\approx7{,}000404$, то есть превосходит порог с
запасом $7{,}6\cdot10^{-5}$ по~$D$.

\emph{Запас.} Для $\ell=103/100$
\[
D(\Gamma)^2-\ell^2\diam^2 V_0
=\frac{5750986852163787427}{147618194625000000000}>0 .
\qedhere
\]
\end{proof}

\begin{remark}[независимая проверка]\label{rem:r7second}
Неравенство $\diam^2V_0<7=D(\Gamma)^2$, а с ним $\chi(\R^7)\le1029$,
получено и вторым путём, не разделяющим с теоремой~\ref{thm:r7ex} ни одной
вычислительной детали: слоёная структура решётки позволяет ограничить радиус
покрытия сверху максимумом явной выпуклой квадратики по вершинам кусков
общего измельчения двух сдвинутых разбиений Вороного \emph{базы}
$E_6^*$ ($167$ кусков, $103$ из них доказано пустыми), причём в
шестимерной базе и без построения семимерной ячейки; максимум равен
$103310189182571717/59047277850000000=1{,}7496\ldots<7/4$, откуда
$\diam^2V_0\le6{,}99847\ldots<7$. Оба вычисления --- целиком в дробях;
подробности --- \cite{Repo}, \cite{Widths}.
\end{remark}

\subsection{Размерность $9$: $\chi(\R^9)\le7203$}\label{ssec:dim9}

\paragraph{Ламинирование.}
Схема та же, что в \S\ref{ssec:dim7}, но базой служит эйзенштейнова раскраска
$E_8/2401$: $\Gamma'=(3+\omega)E_8$, индекс $7^4=2401$, ширина $\sqrt{7/6}$
(предложение~\ref{prop:planar} при радиусе покрытия $E_8$, равном
$\lambda_1/\sqrt2$~\cite{ConwaySloane}). При модуле слоя $m=3$ индекс равен
$2401\cdot3=7203$; высота $t\approx1{,}15$ и сдвиг $c$ найдены
оптимизацией~\cite{Repo}. Проверяется
рациональная точка с $t^2=8264707/6250000$:
% Сгенерировано из audit-data/results/dim9_7203_exact.json (не править вручную)
\begin{equation}\label{eq:gram7203}
\begingroup
\setlength{\arraycolsep}{2.4pt}
G=\begin{pmatrix}
30000&15000&15000&-15000&-15000&15000&15000&15000&3329\\
15000&30000&15000&0&-15000&15000&15000&15000&6390\\
15000&15000&30000&0&0&0&0&0&4411\\
-15000&0&0&30000&0&0&0&-15000&-611\\
-15000&-15000&0&0&30000&-15000&-15000&-15000&-3579\\
15000&15000&0&0&-15000&30000&15000&15000&5606\\
15000&15000&0&0&-15000&15000&30000&15000&4509\\
15000&15000&0&-15000&-15000&15000&15000&30000&8550\\
3329&6390&4411&-611&-3579&5606&4509&8550&17982
\end{pmatrix},\qquad Q=\frac{G}{10^{4}},
\endgroup
\end{equation}
\begin{equation}\label{eq:sub7203}
\begingroup
\setlength{\arraycolsep}{3.0pt}
C=\begin{pmatrix}
-5&-2&-9&-5&-3&-6&-10&-3&0\\
4&3&-4&-3&0&0&-3&0&0\\
0&0&9&5&3&3&8&3&0\\
0&0&-5&-2&0&-3&-4&0&0\\
0&0&-4&-3&0&-1&-3&-1&0\\
0&0&4&3&1&3&3&1&0\\
0&0&0&0&0&0&2&-3&0\\
0&0&0&0&0&0&1&2&0\\
0&0&0&0&0&0&0&0&3
\end{pmatrix},\qquad \det C=-7203 .
\endgroup
\end{equation}

Верхний блок $8\times8$ матрицы $G$ --- решётка $E_8$ в нормировке
$\lambda_1^2=3$ (его определитель равен $(3/2)^8=6561/256$), девятая строка
задаёт слоевой вектор $(c,t)$. Смитова форма матрицы $C$ равна
$\operatorname{diag}(1,1,1,1,1,7,7,7,21)$, то есть
$\Lambda/\Gamma\cong(\Z/7)^3\times\Z/21$: четыре характера $\bmod\,7$ базы
и модуль слоя $3$.

\paragraph{Где было препятствие.}
В отличие от $\R^7$, априорной границы здесь не хватает: для ламинации
$\diam^2V_0\le4R_{\text{базы}}^2+t^2=6+t^2$, что даёт лишь
$R^2\le1{,}8306>7/4$, тогда как порог $D(\Gamma)^2=7$ требует $R^2<7/4$.
Поэтому всё утверждение упирается в один список --- полный набор вершин
девятимерной ячейки. Его удаётся получить после точной унимодулярной замены
базиса: вершин $1\,654\,230$ --- в пятьдесят с лишним раз больше, чем
$30\,368$ в~$\R^7$.

\begin{keybox}
\begin{theorem}[точный сертификат в $\R^9$]\label{thm:r9ex}
Пусть $\Lambda=\Z^9$ со скалярным произведением $x^{\mathsf T}Qy$,
$Q=G/10^4$ из~\eqref{eq:gram7203}, и $\Gamma=C\Z^9$ с $C$
из~\eqref{eq:sub7203}, $|\det C|=7203$. Тогда
\[
D(\Gamma)^2=7,\qquad
\diam^2 V_0=\frac{1119552292542693}{165294140000000},\qquad
d^2=\frac{1157058980000000}{1119552292542693},
\]
то есть $d=1{,}0166127\ldots$, и $\chi\bigl(\R^9,[1,\ell]\bigr)\le7203$ при
всех $1\le\ell<d$; в частности, $\chi(\R^9)\le7203$.
\end{theorem}
\end{keybox}

\begin{proof}[Доказательство --- протокол предложения~\ref{prop:protocol}]
\emph{Ячейка.} Релевантный вектор удовлетворяет $\|v\|^2\le4R^2\le6+t^2$, что
в единицах $G$ даёт границу перебора $73\,223$; из $4504$ перебранных векторов
получаются $376$ пар, то есть $752$ фасеты (потолок в $\R^9$ равен
$2(2^9-1)=1022$: $135$ классов $\Lambda/2\Lambda$ вырождаются и фасеты не
дают). Список фасет в доверенную базу не входит: любой набор целых векторов
решётки задаёт полупространства, содержащие $V_0$, поэтому пропуск фасеты мог
бы лишь увеличить $\diam V_0$ и уменьшить $D(\Gamma)$, то есть ослабить
оценку. Контроль: независимый перебор по классам $\Lambda/2\Lambda$ вовсе без
априорной границы даёт тот же набор из $752$ фасет.

\emph{Вершины.} Как и в теореме~\ref{thm:r7ex}, Qhull лишь \emph{предлагает}
точки, каждая переводится в точную рациональную вершину решением
целочисленной системы по $9$ независимым активным фасетам, а полнота
доказывается точным замыканием по $1$-скелету. Простота \emph{не}
предполагается: из $1\,654\,230$ вершин $1\,649\,640$ простые, $4320$ имеют по
пятнадцать и $270$ --- по шестнадцать активных фасет. Отсюда точное
\[
R^2=\frac{1119552292542693}{661176560000000}=1{,}69327281\ldots<\frac74,
\qquad \diam^2V_0=4R^2 .
\]

\emph{Разделение.} Перебор с подстраховкой $(a+b)^2\le2(a^2+b^2)$, то есть
$\|v\|^2\le2\bigl(D_{*}^2+\diam^2V_0\bigr)=2276611272542693/82647070000000$,
оставляет в окне $280$ векторов ($140$ пар); для каждого предъявляется точный
KKT-сертификат расстояния $D(v)^2=4\dist(v/2,V_0)^2$.

\emph{Минимум} равен $7$ ровно и достигается на $240$ \emph{горизонтальных}
векторах ($120$ пар, нулевая слоевая координата): это в точности пол
$\sqrt{7/3}\,\lambda_1$ предложения~\ref{prop:planar} в нормировке
$\lambda_1(E_8)^2=3$. Ни один слоевой вектор порога не нарушает.

\emph{Запас.} При $\ell=1$
\[
D(\Gamma)^2-\ell^2\diam^2 V_0=\frac{37506687457307}{165294140000000}
=0{,}2269\ldots>0,
\]
а при $\ell=101/100$ он ещё равен
$150036863771988707/1652941400000000000>0$.
\qedhere
\end{proof}

\begin{remark}[что даёт эта строка]\label{rem:r9beats}
Оценка $7203$ вытесняет прежнюю $\chi(\R^9)\le9604$ следствия~\ref{cor:dim9}
сразу по обоим параметрам: цветов меньше, а доказанный отрезок шире
($1{,}0166127\ldots$ против $\sqrt{63/61}=1{,}0162612\ldots$). На оценки в
$\R^{10}$ и выше она не влияет: в них девятимерный блок не участвует.
\end{remark}
